\begin{center}
\begin{tikzpicture}[>=stealth]
 \draw (0,0) rectangle (7.3,0.3);
 \fill (0.15,0.15) circle (0.08) node[left=4pt] {D};
 \draw[->, very thick] (2.4,0.15) -- (2.4,-2.2) node[below] {$F_1$};
 \draw[->, very thick] (6.9,0.15) -- (6.9,0.95) node[above] {$F_2$};
 \draw[|<->|] (0.15,0.75) -- (2.4,0.75) node[midway, above] {$a_1$};
 \draw[|<->|] (2.4,-0.4) -- (6.9,-0.4) node[midway, below] {$d$};
\end{tikzpicture}
\end{center}

Damit sich die Drehmomente um die Drehachse D aufheben, müssen die beiden Kräfte in entgegengesetzte Richtungen zeigen, und die grössere Kraft $F_1 = \kHbO\cdot F_2$ greift näher bei der Achse an, im Abstand $a_1$. Die kleinere Kraft greift im Abstand $a_2 = a_1 + d$ an. Das Hebelgesetz lautet
\begin{align}
 F_1\cdot a_1 &= F_2\cdot\left(a_1 + d\right) \\
 k\,F_2\cdot a_1 &= F_2\cdot a_1 + F_2\cdot d
\end{align}
mit $k = F_1/F_2 = \kHbO$. Aufgelöst:
\begin{align}
 a_1 &= \aHbF = \frac{\dHb}{\kHb-1} = \aHb \approx \result{\aHbP{-2}{2}}
\end{align}
Die kleinere Kraft greift dann im Abstand $a_2 = \bHbP{-2}{2}$ an.
